\chapter{Metodología y planificación}
\label{chap:metodologia_planificacion}

\section{Metodología de desarrollo}
\label{sec:metodologia_desarrollo}

Se decidió emplear una metodología iterativa incremental por lo bien que se adapta tanto a los 
proyectos de este tipo como a la disponibilidad con la que contaban tanto los tutores como el alumno. 

Esta estrategia busca dividir el proyecto en varios bloques conocidos como iteraciones. Cada 
iteración deberá ir aumentando poco a poco las funcionalidades y la complejidad del sistema 
hasta darlo por finalizado. Además, cada iteración, en concreto las dedicadas al desarrollo, 
permite recorrer de forma cíclica las distintas fases de un ciclo de vida tradicional: análisis, 
diseño, implementación y prueba.

Este esquema de trabajo basado en iteraciones se adapta a la perfección con la asistencia al CITIC (Centro de Investigación en 
Tecnologías de la Información y las Comunicaciones) planeada inicialmente. Trabajar de manera presencial en el laboratorio de robótica 
del centro tres veces por semana facilitaba al alumno reunirse frecuentemente con sus tutores y comunicar los avances con ellos.
Aunque inicialmente estaba previsto mantener reuniones periódicas cada 15 días o menos en función de la complejidad de cada tarea, al 
tener el alumno disponibilidad casi total de asistencia, y el tutor estar también disponible, no se llevaron a cabo estas reuniones, sino 
que, en su lugar, se realizó un seguimiento dia a dia en el laboratorio de robótica. Este seguimiento diario permitió detectar, tanto durante el 
planteamiento como durante el desarrollo, fallos y mejoras a implementar, así como la necesidad de retrasar o reducir partes del proyecto.

Otro punto importante para seleccionar esta metodología es la propia naturaleza del proyecto. Al tratarse 
de un trabajo de investigación, especialmente uno que trata sobre aprendizaje automático como este, presenta 
numerosos riesgos (tratados más adelante en la sección \ref{sec:planificacion_inicial}, página \pageref{sec:planificacion_inicial}) 
que podrían dificultar o impedir la realización de parte del trabajo. Este enfoque permite subdividir 
las tareas complejas en unas más sencillas y asequibles. 

Destacar que, gracias a aplicar en cada iteración este ciclo de vida mencionado anteriormente, permitió ir 
obteniendo resultados parciales funcionales y evaluables a lo largo de las etapas antes de finalizar el 
proyecto. Esto fue especialmente útil para el caso de la cancelación de la simulación con MuJoCo y de los problemas 
ocurridos con el entorno Optitrack, como se verá más adelante.

Como último aspecto a favor de esta metodología, cabe destacar que el proyecto implicó el trabajo con numerosas tecnologías 
y conocimientos técnicos que, en primer momento, resultaban desconocidos para el alumno. La puesta a punto del sistema Optitrack, 
del que se disponía del hardware necesario pero cuyo software para la transferencia de datos entre equipos no se había utilizado 
previamente, así como el uso del robot Go2 mediante su \gls{sdk} de Python, también desconocido hasta ese momento, requirieron 
un período de adaptación. A ello se sumó la necesidad de familiarizarse con las distintas librerías empleadas para interconectar 
ambos sistemas y para entrenar el modelo, así como con fundamentos teóricos que no se habían visto con anterioridad, como los 
modelos de mundo y de utilidad, o el aprendizaje automático, tratado de forma introductoria a lo largo de la carrera.

Al subdividir el problema global en tareas concretas enfocadas en tecnologías individuales permiten al alumno familiarizarse 
con cada una de ellas e ir conectándolas a lo largo del proyecto. Además, la cercanía con los tutores gracias a la presencialidad 
facilita la resolución de posibles dudas o problemas que se vayan presentando durante el desarrollo.



\subsection{Iteraciones de la metodología en el proyecto}
\label{sec:iteraciones_metodologia}

Tal como se indicó en el punto anterior, el desarrollo del proyecto se llevó a cabo siguiendo una metodología 
iterativa incremental, con el fin de dividir el trabajo en iteraciones que daban como resultado entregables parciales. 
En cada una se recorrieron de forma cíclica las distintas fases de ciclo de vida adaptadas para el correcto 
desenvolvimiento de este trabajo:

\begin{itemize}
    \item \textbf{Análisis}:
    Durante esta fase se definieron los objetivos y el alcance previsto en cada iteración, determinando qué funcionalidades 
    se implementarían a lo largo de esta. Se trataron los requisitos técnicos necesarios para su realización así como 
    las especificaciones, y se analizaron los posibles riesgos y sus contramedidas. Además, durante esta fase se tuvieron 
    en cuenta retrasos o problemas surgidos durante iteraciones anteriores permitiendo buscar soluciones para mitigarlos o 
    incluso replantear el enfoque en caso de no obtener resultados satisfactorios.

    \item \textbf{Configuración del entorno y formación en las tecnologías}:
    En esta fase, tras haber decidido el objetivo de la iteración, se preparó el entorno y las herramientas necesarias para 
    llevar a cabo la tarea. En base a este objetivo, se realizaron tareas como: configuración del robot Go2 y su conectividad, 
    puesta a punto del sistema de cámaras Optitrack, transferencia de datos entre ordenadores y la creación del entorno de simulación 
    en MuJoCo.
    Asimismo, esta fase permitió al alumno familiarizarse con las nuevas tecnologías y términos necesarios para abordar el desarrollo 
    de cada funcionalidad.

    \item \textbf{Desarrollo}:
    Una vez completadas las fases previas, se procedió a implementar las funcionalidades correspondientes a cada iteración. Los principales objetivos a 
    desarrollar fueron la creación de los scripts para recopilación, almacenamiento y procesamiento de los datos sobre las observaciones que se usarían 
    para entrenar los modelos. Además de implementar los que se utilizarían para controlar el robot, para el entrenamiento de los distintos modelos 
    deliberativos y para obtener las métricas y crear las gráficas con las que analizar los resultados de las experimentaciones realizadas. El conjunto de 
    todas estas funcionalidades terminaron llevando a la creación de un sistema de toma de decisiones basado en un modelo de mundo y de utilidad que 
    permite al robot moverse y cumplir objetivos.

    \item \textbf{Pruebas}:
    Al final de cada iteración se realizó una fase de pruebas para verificar si los objetivos definidos durante la fase de análisis se 
    cumplieron de manera satisfactoria. Estas verificaciones pasan desde pruebas básicas, como verificar la conexión entre Motive y 
    NatNet, hasta algunas más complejas orientadas a los resultados, como la eficacia de los distintos modelos entrenados. En caso de no 
    cumplirse los estándares esperados o de detectar problemas en la realización, estos se anotarán y se analizarían durante la fase 
    correspondiente de la siguiente iteración para realizar las modificaciones pertinentes.

\end{itemize}

Como conclusión, esta metodología permitió obtener resultados parciales evaluables según avanzaba el proyecto y corregir el rumbo 
cuando surgieron situaciones o problemas que dificultaron el avance del trabajo (como se puede observar en la sección \ref{sec:desviaciones_planificacion}, 
página \pageref{sec:desviaciones_planificacion}).

\subsection{Herramientas de apoyo a la metodología}
\label{sec:herramientas_apoyo_metodologia}

A continuación, se describen las herramientas empleadas a lo largo del proyecto, las cuales facilitaron tanto su organización y desarrollo 
como la correcta aplicación de la metodología previamente expuesta.

\begin{itemize}
    \item \textbf{GitHub}\cite{GithubWeb}: GitHub es la plataforma de gestión de versiones, seleccionada para alojar y gestionar el código fuente de este 
    proyecto. Entre sus ventajas destacan la posibilidad de mantener un historial de los avances, recuperar estados anteriores en caso de 
    necesidad y facilitar la comunicación con los demás miembros del proyecto. Además fue especialmente relevante para la metodología utilizada, 
    facilitando la estructuración de las distintas iteraciones y su posterior consulta a modo de índice al marcar la finalización de cada una 
    con un \gls{commit} de los avances realizados.  Por norma general en la gestión de proyectos, suelen utilizarse los branches, ramas distintas a la 
    rama “Master” dedicadas a una tarea en específico y que tras la finalización de esta se fusionan. Sin embargo, en este caso al tratarse de un proyecto 
    personal no se siguió esta norma.

    \item \textbf{Visual Studio Code}\cite{VSCodeWeb}: VS Code es un \Gls{ide} de código abierto ampliamente utilizado en desarrollo de software. 
    Fue el seleccionado para llevar a cabo la implementación del proyecto por características como el soporte de depuración, ejecución de scripts 
    o la integración con GitHub.

    \item \textbf{Overleaf}\cite{OverleafWeb}: Overleaf es una plataforma en línea utilizada para la creación de documentos \LaTeX\cite{LatexWeb}.
    Dispone de una versión gratuita, que fue empleada para redactar esta memoria. Gracias a esta herramienta, fue posible mantener una versión 
    siempre actualizada del documento, lo que facilitó tanto la colaboración con los tutores como la revisión continua del texto.

\end{itemize}


\section{Planificación y seguimiento}
\label{sec:planificacion_seguimiento}

\subsection{Planificación inicial}
\label{sec:planificacion_inicial}

La planificación inicial comenzó con una reunión con los tutores para acordar la temática y los objetivos del proyecto, creando así el 
anteproyecto para este TFG. En esta reunión se definieron los objetivos iniciales, tanto generales como específicos, así como la metodología 
y estrategia diseñadas para alcanzarlos. Posteriormente, se procedió a realizar un planteamiento inicial sobre cómo abordar el trabajo, 
realizando un análisis de las necesidades y la viabilidad, acotando el alcance y los requisitos funcionales, identificando los principales 
riesgos y limitaciones, y especificando las tecnologías necesarias.

\begin{itemize}
    \item \textbf{Análisis de necesidades y viabilidad}

    El objetivo principal del proyecto es disponer de un sistema que permita aprender modelos de mundo y de utilidad a partir de las experiencias del 
    robot, conseguidas mediante la interacción de este con el mundo que lo rodea, para desarrollar un modelo deliberativo capaz de descubrir objetivos 
    y conseguir alcanzarlos en un entorno real. 

    En el caso concreto de este proyecto, el robot debe descubrir que tiene que aprender a acercarse a un determinado objetivo, y posteriormente, 
    realizar el entrenamiento con los datos recopilados para alcanzar ese objetivo de manera consistente.

    Se tuvieron en cuenta dos factores a la hora de evaluar la viabilidad:

    \begin{itemize}
        \item Viabilidad técnica: contar con los elementos necesarios para construir el sistema (robot Go2 y sistema Optitrack), así como las 
        herramientas de desarrollo y los conocimientos requeridos para llevarlo a cabo (Python, Motive y los \Gls{sdk}s). Además, se dispuso del 
        espacio de trabajo adecuado para la puesta a punto y el uso de estos recursos en el laboratorio de robótica del CITIC.

        \item Viabilidad temporal: se contaba con la dedicación completa del alumno (en términos de tiempo) y con el seguimiento periodico de 
        los tutores, junto con la aplicación de una metodología iterativa incremental, que permitió desarrollar el proyecto y corregir 
        desviaciones respecto a la planificación inicial.

    \end{itemize}

    Tras tener estos puntos en cuenta, se concluyó que se trataba de un proyecto viable en el tiempo previsto.

    \item \textbf{Análisis de alcance del proyecto}

    El alcance inicial del proyecto se planteó en dos líneas:

    \begin{itemize}
        \item Desarrollo en entorno simulado (MuJoCo): Creación del entorno, integración de este con el modelo \Gls{xml} del robot y 
        entrenamiento y pruebas de los modelos desarrollados. Funciona como paso previo a trasladar el sistema a un robot real, a modo de 
        testeo de sus funcionalidades y viabilidad.
        
        \item Desarrollo en entorno real: Configuración del entorno Optitrack, transmisión de datos en streaming, control de robot Go2 y 
        entrenamiento y validación de los modelos. Se trata del paso final del proyecto y sirve como prueba de la idea central del trabajo.

    \end{itemize}

    Se ideó siguiendo un planteamiento incremental, priorizando tareas “sencillas”, como la configuración del sistema de cámaras o la 
    transferencia de datos, partes esenciales para el posterior entrenamiento de los modelos. Las funciones no esenciales, como implementación 
    de locomoción a bajo nivel, quedaron descartadas debido a la sobrecomplejidad que suponían.

    \item \textbf{Análisis de requisitos funcionales}

    Los requisitos funcionales describen las funciones y los comportamientos que el sistema debe ser capaz de realizar para cumplir los objetivos 
    especificados. 
    
    Para este proyecto se definieron los siguientes:
    
    \begin{itemize}
    
        \item \textbf{Captura de movimiento del robot y del objetivo}:
        El sistema deberá ser capaz de obtener la posición del robot y del objetivo de forma correcta. Se usará el sistema de captura Optitrack para la 
        obtención de estos datos.
    
        \item \textbf{Transmisión de datos}:
        El sistema deberá soportar la transmisión de datos desde el Motive al ordenador donde se ejecutan los scripts de control y entrenamiento. La 
        transferencia deberá ser estable, garantizando que las observaciones recibidas sean correctas para su posterior uso en entrenamiento.

        \item \textbf{Identificación de elementos del entorno}:
        El sistema deberá ser capaz de diferenciar los datos correspondientes al robot y al objetivo. Con esto, se construirán las observaciones con las 
        que se representará el entorno en cada instante.

        \item \textbf{Control de movimiento del robot}:
        El sistema deberá ser capaz de enviar instrucciones al robot Unitree Go2 con las cuales realizará las acciones designadas previamente. Estas instrucciones 
        corresponden con funciones de alto nivel del \Gls{sdk} del robot y permitirán el movimiento de este a través del entorno.

        \item \textbf{Generación y almacenamiento de datos de entrenamiento}:
        El sistema deberá permitir almacenar los datos asociados a las experiencias de interactuar con el entorno de manera adecuada para su posterior uso en 
        entrenamiento. Para el modelo de mundo deberá guardar estos datos en una lista de tuplas con la forma [\( Obs_t \), Accion, \( Obs_{t+1} \)], mientras que para el modelo de 
        utilidad deberán ser de la forma [\( Obs_{t+1} \), Utilidad] (esto se explicará con más detalle en el apartado de desarrollo, sección \ref{sec:desviaciones_planificacion}, 
        página \pageref{sec:desviaciones_planificacion}).

        \item \textbf{Entrenamiento del modelo de mundo}:
        El sistema deberá ser capaz de entrenar un modelo de mundo que permita predecir el estado siguiente del entorno a partir del estado actual y de la acción a 
        ejecutar por el robot.

        \item \textbf{Entrenamiento del modelo de utilidad}:
        El sistema deberá permitir entrenar un modelo de utilidad con el cual predecir el valor de efectividad que tendrá un determinado estado en función del objetivo 
        que se desea alcanzar.

        \item \textbf{Evaluación de los modelos entrenados}:
        El sistema deberá incluir métodos para validar y evaluar los modelos entrenados, mostrar gráficas y métricas adecuadas para comprobar si las predicciones 
        son suficientemente correctas.

        \item \textbf{Realizar pruebas en entorno real}:
        El sistema deberá ser capaz de ejecutar pruebas reales con el robot con el fin de comprobar su comportamiento en situaciones de uso reales.
        
    \end{itemize}

    A lo largo del proyecto se fueron revisando estos requisitos permitiendo realizar ajustes en las distintas fases y su alcance.

    \item \textbf{Análisis de requisitos no funcionales}

    Los requisitos no funcionales describen las propiedades de calidad y las restricciones que debe cumplir el sistema durante su funcionamiento.
    
    Para este proyecto se definieron los siguientes:

    \begin{itemize}
        \item \textbf{Precisión en la obtención de la posición}:
        La captura de posición del robot y del objetivo deben ser lo suficientemente precisos tanto para que los datos puedan ser utilizados para crear modelos 
        funcionales como para la validación de estos mismos.

        \item \textbf{Transmisión de datos fiable}:
        La transmisión de datos entre Motive y el ordenador receptor deberá ser fiable, con el fin de minimizar posibles fallos o pérdidas en la conexión. 
        Es especialmente relevante la consistencia temporal con las acciones para obtener correctamente los datos de entrenamiento.

        \item \textbf{Seguridad durante las pruebas en entorno real}: 
        El sistema deberá utilizarse siempre en un entorno seguro durante las pruebas con el robot real. Será necesario establecer los límites de movimiento del 
        robot y despejar el entorno para minimizar el riesgo de colisión o caídas durante estas pruebas.

        \item \textbf{Modularidad}:
        El sistema desarrollado deberá estar organizado de manera modular, es decir, separando las distintas fases en códigos independientes (adquisición de 
        datos, entrenamiento, pruebas, control del robot, etc). Esto se realiza con el objetivo de facilitar el mantenimiento y una posible ampliación o modificación 
        posterior, además de adaptarse a la metodología seleccionada, a la planificación inicial y a posibles desviaciones de esta, facilitando su adaptación 
        y justificación.

        \item \textbf{Escalabilidad}:
        Aunque para este proyecto se utilizará únicamente un escenario concreto, el sistema deberá estar diseñado para ser escalable, de forma que pueda 
        adaptarse a nuevos objetivos, conjuntos de datos y modificaciones en los modelos sin necesidad de rehacer por completo el código.

        \item \textbf{Mantenibilidad}:
        El código desarrollado deberá ser claro y estar bien documentado, facilitando la comprensión de este y modificaciones futuras.

        \item \textbf{Eficiencia computacional}:
        El entrenamiento de los modelos deberá presentar un coste computacional razonable para los recursos disponibles y los experimentos deberán poder completarse en 
        un tiempo asequible.

        \item \textbf{Trazabilidad}:  
        El sistema deberá permitir guardar y consultar posteriormente los datos obtenidos de las pruebas realizadas, con el fin de comparar y evaluar los 
        distintos modelos ejecutados durante el desarrollo.
        
    \end{itemize}

    \item \textbf{Análisis de riesgos y limitaciones}

    
    Tras realizar un análisis de riesgos, se identificaron las principales situaciones que podrían comprometer el desarrollo de este 
    proyecto. Para cada riesgo se estimó su probabilidad de ocurrir, así como su impacto sobre el sistema, y se establecieron las medidas 
    de mitigación con la finalidad de prevenir o reducir sus consecuencias.

    Para facilitar la explicación se empleó la siguiente escala:

    \begin{itemize}
        \item Probabilidad: baja (B), media (M), alta (A).
        \item Impacto: baja (B), media (M), alta (A).
    \end{itemize}

    \renewcommand{\arraystretch}{1.1}
    \setlength{\tabcolsep}{3pt}
    \small
    \begin{longtable}{>{\raggedright\arraybackslash}p{0.7cm}|
                    >{\raggedright\arraybackslash}p{2.6cm}|
                    >{\centering\arraybackslash}p{0.8cm}|
                    >{\centering\arraybackslash}p{0.8cm}|
                    >{\raggedright\arraybackslash}p{2.8cm}|
                    >{\raggedright\arraybackslash}p{5.0cm}}
    \caption{Matriz de riesgos del proyecto.}
    \label{tab:matriz_riesgos} \\
    
    \rowcolor{udcpink!25}
    \textbf{ID} & \textbf{Riesgo} & \textbf{Prob.} & \textbf{Imp.} & \textbf{Consecuencias} & \textbf{Medidas de mitigación} \\
    \endfirsthead

    \multicolumn{6}{c}{\tablename\ \thetable{} -- {\small \textit{(vén da páxina anterior)}}} \\
    \rowcolor{udcpink!25}
    \textbf{ID} & \textbf{Riesgo} & \textbf{Prob.} & \textbf{Imp.} & \textbf{Consecuencias} & \textbf{Medidas de mitigación} \\
    \endhead

    \multicolumn{6}{c}{\dotfill{\small \textit{(continúa na páxina seguinte)}}\dotfill} \\
    \endfoot

    
    \endlastfoot

    R1 &
    Problemas en entrenamiento en simulación (falta de funciones de alto nivel, integración con MuJoCo) &
    A & A &
    Retraso importante o eliminación de la parte de simulación &
    Priorizar un modelo simple de locomoción, reutilizar entornos ya existentes, establecer un criterio de parada en caso de consumir demasiado tiempo en esta fase, descartar el entrenamiento en simulación y trabajar únicamente con el robot real. \\
    \hline

    R2 &
    Ausencia de librerías para el control del robot en simulación &
    M & A &
    Imposibilidad de lograr locomoción en el robot en simulación, alto coste de desarrollo &
    Intentar crear funciones de alto nivel, definir acciones sencillas, realizar pruebas progresivas en el desarrollo. \\
    \hline

    R3 &
    Problemas de integración con Optitrack y Motive/NatNet &
    A & A &
    Retraso en la obtención de datos y en el entrenamiento de los modelos. &
    Designar tiempo previamente para la puesta a punto de las tecnologías, pruebas de los servicios de transferencia y streaming con scripts acotados, aislar posibles problemas (ordenador de streaming o de entrenamiento), contactar con el servicio técnico para solucionar el problema. \\
    \hline

    R4 &
    Problemas con el seguimiento de los marcadores o poca fidelidad del tracking &
    M & A &
    Datos no precisos o incompletos, modelos mal entrenados &
    Reubicar cámaras, calibración periódica, definir volumen de captura en el entorno, modificar cantidad y posición de marcadores, filtrar errores. \\
    \hline

    R5 &
    Desincronización en la transmisión de datos &
    M & A &
    Conjuntos de datos $(o_t,a_t,o_{t+1})$ incorrectos, modelos mal entrenados &
    Pruebas de consistencia, timestamping en ambos dispositivos, realizar pruebas sencillas para identificar posible fallo, fijar frecuencia de muestreo de cámaras. \\
    \hline

    R6 &
    Problemas con los datos para entrenamiento (datos insuficientes o poca diversidad) &
    M & M &
    Sobreajuste, mala generalización, modelos mal entrenados &
    Validación con modelos parciales, criterios mínimos de cantidad y calidad, realizar pruebas con distintas posiciones de robot/objetivo, incrementar el muestreo, reutilizar conjuntos de datos de modelos previos. \\
    \hline

    R7 &
    Problemas en el entrenamiento (arquitectura de la red de neuronas, hiperparámetros, no convergencia de los modelos, etc) &
    M & M &
    Retrasos entrenando los modelos, resultados no concluyentes &
    Comenzar con modelos simples e ir aumentando progresivamente, normalizar entradas y salidas, early stopping, métricas eficientes, ajustar los hiperparámetros. \\
    \hline

    R8 &
    Limitaciones computacionales (tiempo de entrenamiento elevado) &
    B & B &
    Retrasos entrenando modelos &
    Reducir tamaño de los modelos, frecuencia y batch ajustados, uso de checkpoints, ejecutar entrenamientos nocturnos, emplear mejores equipos si están disponibles. \\
    \hline

    R9 &
    Cambios en el alcance inicial por dificultades técnicas &
    M & M &
    Desviaciones de objetivos y entregables incompletos &
    Definir entregables mínimos, revisiones periódicas con tutores, justificación en caso de desviaciones en la planificación. \\
    \hline

    R10 &
    Riesgo de daño del material durante las pruebas del robot (caidas, golpes, etc) &
    M & A &
    Daño del robot, interrupción del proyecto &
    Entorno de pruebas despejado y seguro, limite de velocidad para el robot, delimitación clara de los límites seguros, programar parada de emergencia en los scripts de compilación de datos y de validación de entrenamiento, supervisión de las ejecuciones. \\
    \hline

    R11 &
    Problemas con el software (fallos de configuración, errores en el código) &
    M & M &
    Retraso y pérdida de resultados &
    Control de versiones mediante Github, pruebas sencillas en los componentes claves (Motive y NatNet, SDK del robot). \\
    \hline

    R12 &
    Dependencia de disponibilidad del laboratorio o de los recursos (robot, cámaras) &
    B & B &
    Tiempo para realizar pruebas reducido y retrasos en la validación de los modelos &
    Planificar sesiones con antelación, prioridad al trabajo en remoto, priorizar pruebas críticas, organización del trabajo con otros usuarios del laboratorio. \\

    \end{longtable}
    \normalsize

    En resumen, de los riesgos recogidos en la tabla anterior, los más destacables debido a su impacto y probabilidad son:

    \begin{itemize}
        \item R1, R3 y R10, principalmente por su capacidad de retrasar fases completas del proyecto o incluso imposibilitar su finalización.
        \item R2 y R5, debido a que repercuten directamente sobre el entrenamiento del sistema.
    \end{itemize}

    Con el fin de paliar y controlar los posibles problemas que pudieran surgir durante el desarrollo se empleó un seguimiento iterativo: al 
    finalizar cada fase se revisaron los riesgos ocurridos, se adoptaron las medidas de mitigación planteadas y de ser necesario se modificó el 
    alcance y la planificación de ciertas fases para adaptarse a estos inconvenientes.

    \item \textbf{Especificación de tecnologías a utilizar}

    Para llevar a cabo el desarrollo de este trabajo se seleccionaron las siguientes tecnologías (explicadas en mayor detalle 
    en la sección \ref{sec:tecnologias_empleadas}, página \pageref{sec:tecnologias_empleadas}):
    
    \begin{itemize}
        \item Entorno Python: el propio lenguaje y las librerías necesarias para trabajar con los datos y el entrenamiento (Numpy, Tensorflow, etc).
        \item Entorno Optitrack: Sistema de cámaras Optitrack y el software necesario para interactuar con el (Motive y NatNet \Gls{sdk}).
        \item Entorno del robot: Robot Unitree Go2 y su \Gls{sdk} para controlar el movimiento.
        \item Entorno de simulación: MuJoCo y modelo \Gls{xml} del Go2, modificado para añadir el objetivo en el mundo.
    \end{itemize}

    La utilización de todas estas tecnologías de la forma adecuada permite el éxito de este proyecto.

    \item \textbf{Planificación inicial}

    Una vez finalizado todo lo anterior, se procedió a realizar una planificación detallada del proyecto. Para esto se creó un diagrama de Gantt (figura \ref{fig:diagrama_gantt}), en el cual se 
    estructuraron las tareas en fases y se definieron los tiempos estimados para cada una de ellas, así como los responsables de cada una. Gracias a esta planificación, 
    se logró un control adecuado del proyecto en tiempo y recursos, y facilitó solventar las desviaciones ocurridas a lo largo del desarrollo.

    \begin{figure}[hp!]
        \centering
        \includegraphics[width=\textwidth]{imaxes/diagramaGanttTFG.png}
        \caption{Diagrama de Gantt del proyecto.}
        \label{fig:diagrama_gantt}
    \end{figure}

\end{itemize}


\subsection{Recursos humanos}
\label{sec:recursos_humanos}

El desarrollo de este proyecto pudo llevarse a cabo gracias a la colaboración del siguiente equipo:

\begin{itemize}
    \item El alumno cumple la función de desarrollador principal del trabajo, ocupándose de la planificación y el desarrollo de cada una de las 
    iteraciones. Destacando tareas como la configuración de los distintos entornos, la implementación del código para las funcionalidades o la 
    realización de las pruebas de validación necesarias.

    \item Los dos tutores desempeñaron la tarea de supervisar y orientar el desarrollo del proyecto. Su función principal consistio en 
    proporcionar feedback a lo largo de las distintas iteraciones y solucionar dudas y problemas técnicos que surgieron a lo largo del 
    desarrollo del proyecto, así como tomar decisiones sobre el rumbo de este para asegurar la finalización del proyecto.

\end{itemize}

\subsection{Recursos técnicos}
\label{sec:recursos_tecnicos}

Para llevar a cabo el desarrollo de este proyecto fueron necesarios una serie de recursos, tanto de hardware como de software, que hicieron 
posible la implementación y posterior prueba de los distintos modelos desarrollados. A continuación citaremos estos recursos junto a una breve 
descripción: 

\begin{itemize}
    \item \textbf{Hardware}:
    
    \begin{itemize}
        \item Robot Unitree Go2. 
        \item Sistema de cámaras Optitrack (8 cámaras Primex 13W), con el que se obtendrán la posición del robot y del objetivo.
        \item Rack y elementos de soporte y cableado, estructura necesaria para sostener la infraestructura de las cámaras.
        \item Marcadores Optitrack (6 en total) que utilizaran las cámaras para identificar los objetos y obtener su posición.
        \item Intel NUC, ordenador que se utilizará para controlar las cámaras.
        \item Switch para conectar la red de cámaras y el Intel NUC.
        \item Cable Ethernet, el cual se empleará para conectar el ordenador con el robot y deberá tener la longitud suficiente para 
        permitirle moverse por el espacio (en este caso se utilizó un cable de 15m).
        \item Ordenador personal, con potencia suficiente para llevar a cabo el entrenamiento de los modelos de aprendizaje automático.
    \end{itemize}
    \item \textbf{Software}:

    \begin{itemize}
        \item \Gls{sdk} de Unitree, que contiene las funciones de control para el robot. 
        \item Motive, programa desarrollado por Optitrack para controlar las cámaras y obtener los datos de posición.
        \item \Gls{sdk} de NatNet, biblioteca que permite la transferencia de datos vía UDP entre Motive y el ordenador personal.
        \item Python, el lenguaje de programación utilizado para crear los scripts necesarios para el proyecto, junto con las distintas 
        librerías necesarias para esto (mencionadas en la sección \ref{sec:entorno_python}, página \pageref{sec:entorno_python}).
        \item MuJoCo, el motor de físicas empleado para la simulación.
    \end{itemize}
    
\end{itemize}

La combinación de todas estas tecnologías permitió el desarrollo de este trabajo, desde la obtención de los datos necesarios para entrenar 
los modelos de mundo y de utilidad, hasta el posterior entrenamiento y evaluación de los mismos.

\subsection{Costes}
\label{sec:costes}

Para explicar los costes asociados al desarrollo de este proyecto, se dividirán en dos grupos, los referentes al personal y los referentes 
a los materiales necesarios para llevarlo a cabo:

\begin{itemize}
    \item \textbf{Coste de personal}:
    
    \begin{itemize}
        \item Santiago López Silva, como desarrollador principal. Al tratarse del desarrollo de un TFG no obtuvo remuneración por esta tarea, 
        sin embargo, suponiendo un puesto como ingeniero informático junior (recién egresado y con menos de un año de experiencia), el salario 
        medio se encuentra alrededor de los 20.000€ en España (según el portal de empleo 
        Jobted\footnote{\url{https://www.jobted.es/salario/ingeniero-inform\%C3\%A1tico}}).
        \item Martín Naya Varela, con el cargo de ayudante doctor (A-LOSU) para la Universidad de La Coruña, con un salario anual de 
        33.092,08€ (según los datos aportados por la pagina oficial de la UDC\footnotemark).
        \item Alejandro Romero Montero, también con el cargo de ayudante doctor (A-LOSU) para la Universidad de la Coruña, con un salario anual 
        de 33.092,08€ (según los datos aportados por la pagina oficial de la UDC\footnotemark[\value{footnote}]).
        \footnotetext{\url{https://www.udc.es/export/sites/udc/goberno/_galeria_down/xerencia/servizos/retribucions_seguridade_social_e_accion_social/taboas_retributivas/taboas2026/pdi-lab.pdf_2063069294.pdf}}
    \end{itemize}
    
    \item \textbf{Coste material}:
    
    \begin{itemize}
        \item Robot Unitree Go2 EDU PLUS (14.520€).
        \item Sistema de cámaras Optitrack con el software Motive (8 cámaras x 1.500€ y el software 1.000€).
        \item Estructura de soporte para el sistema optitrack, el rack y los soportes para las cámaras (7.  000€).
        \item Intel NUC (500€).
    \end{itemize}

\end{itemize}

Calculando en función de las horas necesarias para realizar este proyecto, da como resultado un coste de 7.060,32€, como se puede ver reflejado en la 
tabla \ref{tab:coste_personal}. Cabe destacar que el material necesario para el proyecto fue proporcionado por el CITIC.

\begin{table}[H]
    \centering
    \small
    \setlength{\tabcolsep}{4pt}
    \renewcommand{\arraystretch}{1.15}
    \rowcolors{2}{white}{udcgray!25}
    \begin{tabularx}{\textwidth}{
    >{\centering\arraybackslash}X|
    >{\centering\arraybackslash}X|
    >{\centering\arraybackslash}p{1.7cm}|
    >{\centering\arraybackslash}p{1.4cm}|
    >{\centering\arraybackslash}p{2cm}}
    \rowcolor{udcpink!25}
    \textbf{Integrante} & \textbf{Rol} & \textbf{€/hora} & \textbf{Horas} & \textbf{Coste} \\
    Santiago López Silva & Desarrollador & 11,63€ & 450 & 5.232,56€ \\
    Martín Naya Varela & Tutor & 19,24€ & 60 & 1.154,37€ \\
    Alejandro Romero Montero & Tutor & 19,24€ & 35 & 673,39€ \\
    \end{tabularx}
    \caption{Coste estimado de personal del proyecto.}
    \label{tab:coste_personal}
\end{table}

\subsection{Desviaciones en la planificación}
\label{sec:desviaciones_planificacion}

A lo largo del desarrollo de este proyecto se experimentaron varios imprevistos que impidieron, en mayor o menor medida, el cumplimiento de la 
planificación inicial. En concreto, se produjeron dos desviaciones relevantes, asociadas con la integración y puesta a punto del entorno 
de simulación y el sistema de captura de movimiento. Estos imprevistos supusieron una reorganización de las prioridades y de los tiempos 
establecidos en un inicio. A continuación, se explicaran sus detalles:

\begin{itemize}
    \item Imposibilidad de realizar el entrenamiento en un entorno simulado:
    Inicialmente se optó por realizar el entrenamiento del modelo deliberativo en un entorno simulado (MuJoCo) como paso previo a su 
    transferencia al robot real, con el fin de probar la eficiencia y la viabilidad del sistema. Sin embargo, durante la implementación de 
    este se encontró un problema que impidió su desarrollo: el control del movimiento del robot en simulación no disponía de las funciones de 
    alto nivel equivalentes a las que se utilizaron posteriormente en el robot real. Esto suponía que para entrenar el robot en este entorno 
    era necesario trabajar a bajo nivel, es decir, operar directamente con acciones sobre los actuadores, controlando el torque aplicado a cada 
    uno de los motores necesarios para la estabilidad y la generación de la marcha.

    Este enfoque acarreaba una complejidad añadida que suponía una desviación importante en el objetivo principal del trabajo (el aprendizaje 
    de modelos de mundo y utilidad), ya que implicaba invertir una gran parte del tiempo previsto en la creación de un controlador de 
    movimiento y en definir un espacio de acciones mucho más complejo. Para solucionar este problema e intentar sacar adelante el entrenamiento de 
    los modelos de mundo y de utilidad en este entorno, se intentó aprender un modelo sencillo de movimiento mediante aprendizaje automático que 
    cumpliera con las acciones básicas con las que se movería el robot (avanzar x metros y rotar sobre sí mismo x grados). Sin embargo, el resultado 
    obtenido no fue satisfactorio, el movimiento del robot era muy rudimentario y no conseguía avanzar.
    Debido a estas pruebas, y a que se vio que aprender un controlador adecuado llevaría más tiempo del deseado, se produjo un retraso 
    significativo. Tras hablarlo con los tutores, se tomó la decisión de priorizar el desarrollo en el entorno real, donde sí existían las funciones 
    de alto nivel (proporcionadas en el \Gls{sdk} del robot) necesarias para hacer que el robot pueda desplazarse y girar de forma correcta, estable y robusta, 
    permitiendo así llevar a cabo la creación del sistema y evitando estancarse en una tarea no prevista.

    Esta decisión llevó a la reducción del alcance esperado para esta fase, reasignando el esfuerzo hacia la creación del sistema en el 
    entorno real, permitiendo terminar el proyecto de forma correcta dentro del plazo establecido.
 
    \item Retraso en la configuración del Motive y la transmisión de datos:
    La segunda desviación se produjo durante la fase de puesta a punto del sistema Optitrack. A pesar de contar con el hardware necesario, la 
    configuración del Motive para realizar de manera efectiva la transmisión de datos hacia el otro ordenador (donde se ejecutaban los scripts 
    de entrenamiento), tomó más tiempo del esperado inicialmente. 

    Fue necesario realizar esta comunicación, en lugar de ejecutar todo en el mismo ordenador, entre otras porque Motive solo tiene soporte para Windows 
    y no para Ubuntu. Además, el fabricante recomendaba mantener separado Motive para evitar problemas con la actualización.


    Este retraso se debió principalmente a dos factores: el primero debido al desconocimiento inicial de la tecnología (tanto de Motive como 
    del \Gls{sdk} de Natnet) así como la configuración de red necesaria requirieron más tiempo del previsto para conseguir que todo funcionara 
    de manera correcta. Realizar pruebas para descubrir el modo de conexión correcto para la configuración utilizada (modo Unicast) y 
    configurar las IPs y los puertos de los distintos dispositivos, además de estudiar y modificar algunas funciones del SDK para recibir los 
    paquetes de información (vía UDP) en el ordenador receptor.
    El otro factor desencadenante del retraso en esta parte fue la tardanza del servicio técnico español de Optitrack a la hora de resolver 
    ciertas dudas planteadas al comienzo de la fase de familiarización con la tecnología, teniendo que contactar con el soporte técnico de 
    Reino Unido para resolverlas.

    El resultado de este retraso llevó al bloqueo de las tareas posteriores (como la toma de datos y el entrenamiento de los modelos) hasta 
    ser capaces de establecer un flujo de datos estable que permitiera obtener estos datos de manera correcta.


\end{itemize}

Ambos retrasos reflejan la naturaleza experimental de proyectos de este estilo y la necesidad de replanificación y tomar contramedidas que 
permitan sobrepasar incidentes durante el proyecto. Finalmente se consiguió completar el tema principal del proyecto, demostrando que a pesar 
de tener que prescindir de una parte concebida inicialmente como importante, el proyecto seguía siendo viable y cumplía los puntos 
fundamentales planteados al inicio.

